% concatenated from main.tex
\documentclass[11pt]{amsart}

\usepackage[english]{babel}
\usepackage{txfonts}
\usepackage{amsmath}
\usepackage{amssymb}
\usepackage{mathrsfs}
\usepackage{color}
\usepackage{esint}
\usepackage{enumerate}

\usepackage[colorlinks=true,
  linkcolor=blue,
  citecolor=red,
  urlcolor=magenta,
  backref=false]{hyperref}

\usepackage[T1]{fontenc}
\usepackage{lmodern}    

\usepackage{anysize}

\textwidth=15cm
\textheight=25cm
\oddsidemargin 0.46cm
\evensidemargin 0.46cm
\parindent=13pt

\allowdisplaybreaks
\raggedbottom

\newtheorem{theorem}{Theorem}[section]
\newtheorem{lemma}[theorem]{Lemma}
\newtheorem{corollary}[theorem]{Corollary}
\newtheorem{proposition}[theorem]{Proposition}

\theoremstyle{definition}
\newtheorem{definition}[theorem]{Definition}
\newtheorem{example}[theorem]{Example}
\newtheorem{remark}[theorem]{Remark}

\newtheorem*{problem*}{Problem}
\numberwithin{equation}{section}

\DeclareMathOperator{\h}{\mathbf{h}} % Metric functional
\newcommand{\abs}[1]{\lvert#1\rvert} %  Absolute value
\newcommand{\norm}[1]{\lVert#1\rVert} %  Norm 
\newcommand{\R}{\mathbb{R}}  % The real numbers.
\newcommand{\N}{\mathbb{N}}  % The natural numbers.

% To include the section number in the equation numbering:
\numberwithin{equation}{section}

\begin{document}

\title{A Weak Topology on Metric Spaces}
\begin{abstract}
We use metric functionals to build a weak topology associated 
with $d$-weak convergence on metric spaces. We also show the
existence of a sequence in the space $C[0,1]$ that is unbounded 
in the sup-norm yet converges $d$-weakly. 
This points out an error in an earlier work.
\end{abstract}
\author{Armando W. Gutiérrez}
\author{Olavi Nevanlinna}

\date{}

\maketitle

\section{Main results}

In an earlier work \cite{GNweakconv}, we introduced the notion 
of $d$-weak convergence in arbitrary metric spaces $(X,d)$. We 
restricted our analysis there to sequences with the purpose of
testing and validating our ideas in a simple way. We verified that
$d$-weak convergence behaves as expected in \emph{all} normed 
linear spaces; every \emph{bounded sequence} in a normed linear space 
converges $d$-weakly if and only if it converges in the classical 
weak sense.

Our goal here is to build a topology on metric spaces 
that is connected to our notion of weak convergence. Let us 
recall first the formal definition of $d$-weak convergence. 

\begin{definition}
Let $(X,d)$ be a metric space and $X^\diamondsuit$ the space of 
all metric functionals on $X$. We say that a net of points 
$(a_i)_{i\in I}$ in $X$ converges $d$-weakly to a point $a$ in $X$ 
if for every $\h$ in $X^\diamondsuit$ we have 
$$
    \liminf_{i\in I}\h(a_i) \geq \h(a).
$$
\end{definition}

Our first result is the following. 

\begin{theorem}\label{thm:thm1}
Let $(X,d)$ be a metric space and $X^\diamondsuit$ the space of 
all metric functionals on $X$. Then, there exists a topology on $X$ 
induced by $X^\diamondsuit$, denoted by $\sigma(X,X^\diamondsuit)$,
with the following properties:
\begin{enumerate}
    \item The topology $\sigma(X,X^\diamondsuit)$ is coarser 
    than the metric topology on $X$.
    \item The topology $\sigma(X,X^\diamondsuit)$ on $X$ makes 
    every metric functional lower semi-continuous.
    \item A net $(a_i)_{i\in I}$ in $X$ converges in the 
    topology $\sigma(X,X^\diamondsuit)$ if and only if it 
    converges $d$-weakly.
\end{enumerate}
\end{theorem}

In \cite[Conjecture 1.5]{GNweakconv} we conjectured
that $d$-weakly convergent sequences in normed linear spaces
must be bounded. The conjecture was based
on a claim in \cite[Theorem 1.4]{GNweakconv}. Such a 
claim contains an error. The following theorem shows that
the conjecture is false.

\begin{theorem}\label{thm:conjfalse}
There exists a sequence $(f_n)$ in $C[0,1]$ that
converges to $0$ in the topology 
$\sigma(C[0,1],C[0,1]^\diamondsuit)$ and 
$\norm{f_n}_\infty \to \infty$.
\end{theorem}

\begin{theorem}\label{thm:thm2}
Let $(X,d)$ and $(Y,D)$ be metric spaces. Suppose that
$T:X \to Y$ is a distance-preserving mapping, that is, 
$D(Tp,Tq) = d(p,q)$ for all points $p,q$ in $X$.
Assume that $Y=TX$. Then, the mapping
\[ T: (X, \sigma(X,X^\diamondsuit)) \to (Y, \sigma(Y,Y^\diamondsuit))\]
is a homeomorphism.
\end{theorem}

It is not an overstatement to say that $\sigma(X,X^\diamondsuit)$ 
is a \emph{weak topology} on $X$. Indeed, our topology has a
connection to the notion of weak metrics, which
was introduced by H. Ribeiro \cite{WeakMetric} in 1943. 

\begin{definition}
Let $X$ be a set. A mapping $\delta: X\times X \to [0,+\infty)$
is called a weak metric if we have $\delta(x,x)=0$ for all
$x$ in $X$ and $\delta(x,y)\leq \delta(x,z)+\delta(z,y)$ 
for all $x,y,z$ in $X$.
\end{definition}

Examples of weak metrics are found in \cite{WeakMetED}. 
Some fixed point problems in weak metric spaces 
were investigated in \cite{FixedPWeakM}.
The following theorem reveals a connection between
$d$-weak convergence and weak metrics.

\begin{theorem}\label{thm:thm3}
Let $(X,d)$ be a metric space and $X^\diamondsuit$ the space of 
all metric functionals on $X$. Then, there exists a family of 
weak metrics $\{ \delta_{\h} \mid \h \in X^\diamondsuit \}$ on $X$ 
such that the following properties are equivalent:
\begin{enumerate}
    \item A net $(x_i)_{i\in I}$ in $X$ converges to
    $x$ in $X$ in the topology $\sigma(X,X^\diamondsuit)$.
    \item $\delta_{\h}(x,x_i)\to 0$ for all $\h$ in $X^\diamondsuit$.
\end{enumerate}
\end{theorem}

\section{Proofs}\label{sec:proofs}

Let us begin recalling the concept of a metric functional. 
Let $(X,d)$ be a metric space. Choose a base-point $o$ 
in $X$ and define the set 
\begin{equation}
	X^{\vee} := \{ d(\cdot,w) - d(o,w) \mid w \in X\}.
\end{equation}
Note that each element in $X^{\vee}$ is a 1-Lipschitz mapping
$X\to \R$ vanishing at the point $o$. Let $X^{\diamondsuit}$ 
denote the closure of $X^{\vee}$ in the topology of pointwise
convergence. Each element in $X^{\diamondsuit}$ is called a 
metric functional on $X$ and each element in $X^{\vee}$ is 
called internal. For more details see \cite{GNweakconv}.

We now move on to building our weak topology on $X$. 
For a point $x$ in $X$, a finite subset $F$ 
of $X^\diamondsuit$ and a positive real number $c$, define 
the set
$$
    B(x,F,c):= \bigcap_{\h\in F}
    \{ y\in X \mid\, \h(y) > \h(x) - c\}.
$$

\begin{lemma}
The collection 
$$
\mathscr{B}:= \{ B(x,F,c) \mid x\in X,\, 
F \text{ finite subset of }X^\diamondsuit,\, c > 0\}
$$
is a basis for a topology on $X$.
\end{lemma}

\begin{proof}
First, note that every point $x$ in $X$ is contained in 
$B(x,\{\h_x\},1)$.
Now, suppose that $x$ is a point in the intersection
of two elements of $\mathscr{B}$, say $B(x_1,F_1,c_1)$ 
and $B(x_2,F_2,c_2)$. The following two inequalities hold:
\begin{enumerate}
    \item $\h(x) > \h(x_1) - c_1$, for all $\h$ in $F_1$,
    \item $\h(x) > \h(x_2) - c_2$, for all $\h$ in $F_2$.
\end{enumerate}
Define $F:=F_1 \cup F_2$ and
\begin{align*}
    c &:= \min\left\{\min_{\h\in F_1}\{\h(x)-\h(x_1)+c_1\}, 
    \min_{\h\in F_2}\{\h(x)-\h(x_2)+c_2\}\right\}. \\
\end{align*}
Then, $B(x,F,c)$ is an element in $\mathscr{B}$, contains
the point $x$, and is a subset of the intersection of 
$B(x_1,F_1,c_1)$ and $B(x_2,F_2,c_2)$.
\end{proof}

\begin{definition}
We denote by $\sigma(X,X^\diamondsuit)$ the topology generated 
by the basis $\mathscr{B}$. That is, a subset $A$ of $X$ is 
open in the topology $\sigma(X,X^\diamondsuit)$ if and only 
if for every $x$ in $A$ there exists a basis element $B$ 
in $\mathscr{B}$ such that $x\in B \subset A$.    
\end{definition}

Note that a subset $A$ of $X$ is open in the topology
$\sigma(X,X^\diamondsuit)$ if and only if for every $x$ 
in $A$ there exists a finite subset $F$ of $X^\diamondsuit$ 
and a positive real number $c$ such that $B(x,F,c)$ is a 
subset of $A$. 

One may wonder if different base-points produce different 
topologies. We show next that the choice of the base-point 
is irrelevant.

\begin{proposition}
Let $p$ be a point in $X$ such that $d(o,p)> 0$. 
Let $X_p^\diamondsuit$ denote the space of all metric functionals 
on $X$ with base-point $p$. Then, we have
$$
    \sigma(X,X^\diamondsuit)=\sigma(X,X_p^\diamondsuit).
$$
\end{proposition}
\begin{proof}
For every $x,y$ in $X$ we have 
$$ d(x,y)-d(o,y)-(d(p,y)-d(o,y)) = d(x,y)-d(p,y).$$
Then, $\h - \h(p)$ is in $X_p^\diamondsuit$ for all $\h$ in 
$X^\diamondsuit$. Now, if the points $o$ and $p$ are permuted 
in the previous equality, we conclude that the mapping 
$\h\mapsto \h - \h(p)$ is a bijection from $X^\diamondsuit$ 
onto $X_p^\diamondsuit$.
\end{proof}

\begin{lemma}\label{lem:coarse}
The topology $\sigma(X,X^\diamondsuit)$ is coarser than 
the metric topology on $(X,d)$. 
\end{lemma}

\begin{proof}
Let $A$ be an open set in the topology $\sigma(X,X^\diamondsuit)$
and $x$ a point in $A$. Then, there exists a basis element 
$B(x,F,c)$ which is a subset of $A$. Since every metric 
functional is $1$-Lipschitz, the set 
$U(x,c):=\{y\in X \mid d(x,y) < c\}$ is a subset of 
$B(x,F,c)$. Since $U(x,c)$ is a ball, which is
open in the metric topology, it follows that the set $A$ 
is open in the metric topology.
\end{proof}

Metric functionals are always continuous as mappings 
$X\to\R$ with their corresponding metric topologies. 
If $X$ is equipped with the topology $\sigma(X,X^\diamondsuit)$,
we have the following.

\begin{lemma}\label{lem:lowersc}
The topology $\sigma(X,X^\diamondsuit)$ makes 
every metric functional lower semi-continuous on $X$.  
\end{lemma}

\begin{proof}
Let $\h$ be a metric functional on $X$ and $r$ a real number. 
Suppose that $x$ is a point in $X$ with $\h(x) > r$. Then,
for every $y$ in the basis element $B(x,\{\h\}, \h(x)-r)$
we have $\h(y) > r$.
\end{proof}

\begin{proof}[\bf Proof of Theorem \ref{thm:thm1}]
The first property is shown in Lemma \ref{lem:coarse} and 
the second property in Lemma \ref{lem:lowersc}.

Suppose that a net $(a_i)_{i\in I}$, with $(I,\succcurlyeq)$
a directed index set, converges $d$-weakly to a point $a$ 
in $X$. Let $A$ be an open set in the topology 
$\sigma(X,X^\diamondsuit)$ and assume that the point $a$ 
is in $A$. Then, there exists a positive real number $c$ 
and a finite subset $F$ of $X^\diamondsuit$ such that the 
basis element $B(a,F,c)$ is a subset of $A$.
By $d$-weak convergence, for each metric functional $\h$ 
in $F$ there exists an index $i_0(c,\h)\in I$  such that 
$$
    \h(a_i) > \h(a) - c,
$$
for all $i \succcurlyeq i_0(c,\h)$.
Now, since $F$ is a finite set, there exists an index $i_*\in I$ 
such that $i_* \succcurlyeq i_0(c,\h)$ for all $\h$ in $F$. 
Then, for every $i \succcurlyeq i_*$ we have 
$$
  a_i \in B(a,F,c) \subset A.  
$$
This shows that $(a_i)_{i\in I}$ converges to $a$ in the 
topology $\sigma(X,X^\diamondsuit)$.

Now, suppose that a net $(a_i)_{i\in I}$ converges to a point
$a$ in $X$ in the topology $\sigma(X,X^\diamondsuit)$. 
Let $\h$ be in $X^\diamondsuit$ and $c$ a positive real number. 
Since $B(a,\{\h\},c)$ is open in the topology 
$\sigma(X,X^\diamondsuit)$, there exists an index $i_0$ 
such that 
$$
    a_i \in B(a,\{\h\},c),
$$
for all $i \succcurlyeq i_0$. In other words, we have
$$
    \liminf_{i\in I} \h(a_i) \geq \h(a).
$$
\end{proof}

\begin{proof}[\bf Proof of Theorem \ref{thm:conjfalse}]
Choose a sequence of closed intervals
$A_n\subset[0,1]$ such that
$A_n\cap A_m = \emptyset$ for all $n\neq m$ and $|A_n|\to 0$. 
For each $n\geq 1$, 
let $f_n$ be a non-negative continuous function on $[0,1]$
that has non-zero values on the interval $A_n$
and a maximum value equal to $n$. We have $\norm{f_n}_\infty = n$
for all $n\geq 1$. We show next that for every
$\h$ in $C[0,1]^\diamondsuit$ we have
\[
 \liminf_{n}\h(f_n) \geq 0.
\]

First, we claim that for each internal $\h_w \in C[0,1]^\vee$,
where $\h_w(\cdot)=\norm{\cdot - w}_\infty-\norm{w}_\infty$, the
value $\h_w(f_n)$ is negative for at most one $n$. Indeed, for a
fixed $w\in C[0,1]$ there exists $s\in [0,1]$ such that
$\abs{w(s)}=\norm{w}_\infty$. If $s$ is not in $A_n$, then
$f_n(s)=0$ and 
\[
    \h_w(f_n)=\norm{f_n - w}_\infty - \norm{w}_\infty \geq \abs{w(s)}-\abs{w(s)}=0.
\]
Since $A_n$ and $A_m$ are disjoint for all $n\neq m$, the
value $\h_w(f_n)$ is negative for at most one $n$.

Now, let $\h$ be an arbitrary metric functional on $C[0,1]$.
Then, there exists a net $(w_i)_{i\in I}$ in $C[0,1]$ such that
$\h_{w_i}(f) \to \h(f)$ for all $f$ in $C[0,1]$. Suppose that 
there exist two distinct positive integers $n$ and $m$ such that
$\h(f_n) < 0$ and $\h(f_m) < 0$. Choose $\epsilon >0$ such that
$\h(f_n)< -2\epsilon$ and $\h(f_m)< -2\epsilon$. Since
$\h_{w_i} \to \h$, there exist $i(\epsilon,n)$ and $j(\epsilon,m)$
in the directed set $(I,\succcurlyeq)$ such that 
$\h_{w_i}(f_n) < \h(f_n) + \epsilon$
for all $i\succcurlyeq i(\epsilon,n)$ and 
$\h_{w_i}(f_m) < \h(f_m) + \epsilon$ for all $i\succcurlyeq j(\epsilon,m)$.
Let $i_0$ be the index in $I$ such that $i_0 \succcurlyeq i(\epsilon,n)$ and 
$i_0 \succcurlyeq j(\epsilon,m)$. Then, for every $i\succcurlyeq i_0$ 
we have
\[
 \max\{\h_{w_i}(f_n),\h_{w_i}(f_m)\} < -\epsilon
\]
But according to the previous argument, each internal $\h_{w_i}$ has 
negative value for at most one $f_n$. Therefore, $\h(f_n) < 0$
for at most one $n$. This implies 
\[
 \liminf_{n}\h(f_n) \geq 0 =\h(0).
\]
\end{proof}

\begin{remark}
We listed in \cite{GNweakconv} some normed linear spaces where
$d$-weak convergence of sequences implies boundedness. 
However, the argument we gave for the case $C[0,1]$ contains an 
error, which is revealed by the sequence $f_n$ 
presented above.
In $\ell_1$ and in a normed linear space $X$ with strictly convex 
dual, $d$-weakly convergent sequences are bounded.
\end{remark}

\begin{proof}[\bf Proof of Theorem \ref{thm:thm2}]
Choose $o$ in $X$ to build the space $X^\diamondsuit$
and choose $To$ in $Y$ to build the space $Y^\diamondsuit$. 
Since $T$ is surjective and distance-preserving, for every
point $y$ in $Y$ there exists a point $w$ in $X$ such that
for every $x$ in $X$ we have
\[
    D(Tx,y)-D(To,y) = d(x,w)-d(o,w).
\]
Then, the mapping 
\[ T^\diamondsuit : Y^\diamondsuit \to X^\diamondsuit 
\quad \h \mapsto \h\circ T\]
is a bijection. 

Now, if a net $(x_i)_{i\in I}$ in $X$ converges $d$-weakly
to $x$ in $X$, then for every $\h$ in $Y^\diamondsuit$ we 
have
\[
    \liminf_{i\in I}\,\h(Tx_i) 
    = \liminf_{i\in I}\,T^\diamondsuit\h(x_i)
    \geq T^\diamondsuit\h(x)
    = \h(Tx).
\]
This shows that $(Tx_i)$ converges $D$-weakly to $Tx$ in $Y$.

On the other hand, if $(Tx_i)_{i\in I}$ in $Y$ converges 
$D$-weakly to $Tx$ in $Y$, then for every $\h$ in $X^\diamondsuit$
we have
\[
    \liminf_{i\in I}\,\h(x_i) 
    = \liminf_{i\in I}\,(T^\diamondsuit)^{-1}\h(Tx_i)
    \geq (T^\diamondsuit)^{-1}\h(Tx)
    = \h(x).
\]
This shows that $(x_i)$ converges $d$-weakly to $x$ in $X$.
\end{proof}

\subsection{Weak metrics}
For each $\h$ in $X^\diamondsuit$ define the mapping 
$\delta_{\h} : X\times X \to [0,+\infty)$ by
$$
    (x,y) \mapsto \delta_{\h}(x,y):= \max\{0, \h(x) - \h(y)\}.
$$

\begin{lemma}
For every $\h$ in $X^\diamondsuit$, $\delta_{\h}$ is a weak 
metric on $X$.    
\end{lemma}

\begin{proof}
Recall that for every real numbers
$a$ and $b$ we have
\[
 0 \leq \max\{0,a+b\} \leq \max\{0,a\} + \max\{0,b\}.
\]
\end{proof}

\begin{proof}[\bf Proof of Theorem \ref{thm:thm3}]
A net $(x_i)_{i\in I}$ in $X$ converges to
$x$ in $X$ in the topology $\sigma(X,X^\diamondsuit)$
if and only if for every $\h$ in $X^\diamondsuit$ and
a positive real number $\epsilon$ there exists $i_0$ 
in $I$ such that
\[
    \h(x_i) > \h(x) - \epsilon,
\]
for all $i\succcurlyeq i_0$. This is equivalent to
the property $\delta_{\h}(x,x_i)\to 0$ for all 
$\h$ in $X^\diamondsuit$.
\end{proof}

In general, the topology $\sigma(X,X^\diamondsuit)$ is not 
Hausdorff. For example, if $X$ is the real line equipped
with the metric $d(x,y):=\sqrt{\abs{x-y}})$, we have
$X^\diamondsuit = X^\vee \cup\{\mathbf{0}\}$ and the 
sequence $(a_n)$ with $a_n:=n$ in $X$ converges
in the topology $\sigma(X,X^\diamondsuit)$ to every 
point in $X$.
For normed linear spaces, the topology $\sigma(X,X^\diamondsuit)$
behaves better.

\begin{proposition}[{\cite[Theorem~1.2]{GNweakconv}}]
Let $X$ be a normed linear space. Then, the 
topology $\sigma(X,X^\diamondsuit)$ is Hausdorff.  
\end{proposition}
\begin{proof}
The proof of \cite[Theorem~1.2]{GNweakconv} is valid for nets.    
\end{proof}
  
\bibliographystyle{plain}
\bibliography{ref}
\end{document}
